\chapter{Motion, rounding and arithmetic range}
\MARSChapterHeading{Motion, rounding and arithmetic range}
\label{revision-lab:fixed-point}

Import the development lesson \texttt{012-fixed-point} and choose Continue
in Learn. Its \texttt{fixed-point-ball.asm} is the same complete program supplied
in \texttt{revision/labs/files}; it follows original project
\texttt{012-fixed-point-ball}. Read the main chapter's raw-value interpretation
and reflection rule first. Learn supplies inline predictions, source pointers,
eight ordered Debug observations, ordinary Run and two independent checks.
This supplement explains what those observations mean and provides experiments
beyond the checked route.

Keep the guided source at \texttt{DEMO} zero, \texttt{round\_mode} zero and eight
steps. Use separate copies for nearest-cell rendering and the overflow probe.
Those variants intentionally have different observations. An edit invalidates
execution credit for the previous source revision; Continue returns to the
first check that needs fresh evidence while retaining the prediction history.

\section{Predict the model before drawing it}

Convert the four default raw seeds to pixels and pixels per update. Derive ticks
zero, one and two by hand, retaining raw integers through every addition. For each
tick, compute the display ID in both modes. Explain why the two modes can choose
different IDs while keeping identical positions and velocities.

Trace the vertical crossing from tick four to five. Derive the trial, its excess
past the upper wall, the reflected raw position and the returned velocity. Compare
that result with clamping at the wall. Continue both hypothetical trajectories for
one more tick to show why the lost overshoot matters later.

Repeat for an exact contact: position 1664 and velocity 128. The returned position
is at the upper wall, but the returned velocity must point inward. Then consider
position zero and velocity $-1024$. The accepted initial state is outward-facing;
the first update reflects it inward. State why one reflection suffices under the
actual speed bound, and why that argument would fail for unrestricted velocities.

\section{Inspect a complete tick and an intermediate call}

Break at \texttt{frame\_ready}. At the first pause, inspect the four raw words,
tick and pixel ID. It is tick zero, not the first completed update. The frame must
contain exactly one white word at the reported ID and black elsewhere. A request
of zero still requires this initialization and publication.

Step past the marker to \texttt{frame\_continue}, then follow the horizontal
\texttt{advance\_axis} call. After its return, inspect both results and the
saved-register moves that retain them. During the next, vertical call, memory
still describes the preceding publication. Explain which registers now contain
new horizontal state and which still contain old vertical state.

Continue to the next ready boundary and verify that both axes, tick, ID and
frame agree. Each requested update adds one publication to the initial one.
Do not count tool repaints during a fast Run: the screen can skip intermediate
frames even though the simulator executes each render.

For the guided route, move one line breakpoint at a time and resume to the
next requested stop. Reusing a label means stopping at its next visit; a pause
is before the selected instruction executes. The first \texttt{frame\_ready}
shows tick zero and pixel nine. At the first \texttt{advance\_axis}, arguments
are horizontal raw position 320 and velocity 192. At the next call, the new
horizontal result 512 is already in \texttt{\$s0}; the vertical arguments are
448 and 320. The following ready boundary publishes tick one, position
$(512,768)$ and pixel 26. This mixed register state between calls is expected;
it is not yet a complete published tick.

Continue to \texttt{reflect\_high}. The vertical trial is 2048 before the
reflection instructions execute. The published memory still describes tick
four, $(1088,1728)$, even though \texttt{\$s0} already holds the next horizontal
position 1280. The next ready boundary publishes tick five at $(1280,1536)$
with vertical velocity $-320$ and pixel 53. At the next upper reflection,
the horizontal trial is 1856; the old horizontal position is 1664. Finally,
\texttt{frame\_ready} publishes tick eight at $(1728,576)$, velocity
$(-192,-320)$ and pixel 22. Finish or stop Debug before starting ordinary Run.
It should exit silently. Inspect guards as well as the visible white pixel.

Set mode to one in a copy and inspect ticks seven and eight. In both frames,
the ball is in column seven. At tick seven its raw $x$ is 1664 and velocity is
192; at tick eight they are 1728 and $-192$. Explain which tick precedes contact
and why the same displayed column does not imply the same physical state.
Figure~\ref{fig:revision-fixed-contact} shows those actual pauses.

\begin{figure}[htbp]
\centering
\begin{tikzpicture}[alt={Two actual native eight by eight Bitmap Display observations in nearest mode. Tick seven shows a white cell at column seven, row four, ID 39; raw x is 1664 with velocity 192. Tick eight shows a white cell at column seven, row two, ID 23; raw x is 1728 with velocity minus 192. All other cells are black. Native double outlines mark column seven in each panel. The first state precedes contact; the second follows reflection.}]
\node[anchor=north west,inner sep=0,font=\small\bfseries] (atitle) at (0,0) {1. Tick seven: before contact};
\node[anchor=north west,inner sep=0] (a) at ([yshift=-2mm]atitle.south west) {\MARSRevisionCrop{revision/captures/fixed-point/ball-tick-7.png}{0.31\linewidth}{0.203906250}{0.066666667}{0.203906250}{0.091111111}};
\coordinate (borigin) at ([xshift=15mm]a.north east);
\node[anchor=north west,inner sep=0,font=\small\bfseries] (btitle) at (borigin |- atitle.north west) {2. Tick eight: after reflection};
\node[anchor=north west,inner sep=0] (b) at (btitle.south west |- a.north west) {\MARSRevisionCrop{revision/captures/fixed-point/ball-tick-8.png}{0.31\linewidth}{0.203906250}{0.066666667}{0.203906250}{0.091111111}};
\foreach \panel in {a,b} {
 \begin{scope}[shift={(\panel.south west)},x={($(\panel.south east)-(\panel.south west)$)},y={($(\panel.north west)-(\panel.south west)$)}]
 \draw[black,line width=1.6pt] (.875,0) rectangle (1,1);
 \draw[white,line width=.65pt] (.875,0) rectangle (1,1);
 \end{scope}
}
\end{tikzpicture}
\caption{Actual light-theme Debug pauses of the public program in nearest mode. Both outlined columns are seven, but the model changes from raw $(x,v_x)=(1664,192)$ to $(1728,-192)$. The white cells have IDs 39 and 23. The column outlines and panel labels are native publication annotations.}
\label{fig:revision-fixed-contact}
\end{figure}


\section{Diagnose quantization and reflection faults}

In a disposable copy, remove the half-unit rounding bias. Find the first mismatch
against nearest mode using raw 127, 128 and 129. Restore the bias, then check raw
1664. The displayed wall column is selected before the actual raw wall is reached.
Rendering must not use that display cell to decide whether to reverse velocity.

For another copy, replace a saved raw position with its integer-pixel part after
rendering. Predict how repeated updates lose fractional displacement, then compare
the first incorrect raw tick with the working version. Repair the model value;
changing only a pixel or a caption cannot restore the lost fraction.

Change the upper contact comparison from inclusive to strict. An exact contact
may still draw at the right place, so inspect the returned velocity. Test the
lower contact separately. A frame-only test can miss this error until the next
update. Restore both contact conventions and test an overshoot as well as equality.

Finally call \texttt{render\_ball} with \texttt{spare\_frame} as destination.
It should return the same ID without writing the normal display frame. This
separates a helper's pointer contract from main's fixed display choice. Run both
original checks against your source from the public repository root:
\begin{verbatim}
java -jar /path/to/cli.jar grade \
  packages/012-fixed-point-ball/checks/kernels.kts \
  /path/to/your/fixed-point-ball.asm
java -jar /path/to/cli.jar grade \
  packages/012-fixed-point-ball/checks/program.kts \
  /path/to/your/fixed-point-ball.asm
\end{verbatim}

\section{Compare a wrapped sum with an exact wider sum}

Set \texttt{DEMO} to one in a copy. Break at \texttt{overflow\_ready}, before the
exit service replaces \texttt{\$v0}. Inspect the full 64-bit values of
\texttt{\$v0} and \texttt{\$v1}. Derive the exact sum of raw hexadecimal
\texttt{7FFFFF00} and \texttt{00000200}, then explain the sign extension in the
wrapped word result. Both results use the same scale; only the arithmetic width
and retained range differ. Compare the actual native register observation in
Figure~\ref{fig:revision-fixed-overflow}.

\begin{figure}[htbp]
  \centering
  \begin{tikzpicture}[alt={An actual native Debug register observation before the exit service. Decimal v0 is minus 2147483392, the sign-extended wrapped word sum. Decimal v1 is 2147483904, the exact positive doubleword sum. A native publication rectangle encloses these two adjacent register rows.}]
    \node[inner sep=0] (values) {\MARSRevisionCrop{revision/captures/fixed-point/ball-overflow.png}{0.45\linewidth}{0.243750000}{0.313333333}{0.601562500}{0.642222222}};
    \begin{scope}[shift={(values.south west)},x={($(values.south east)-(values.south west)$)},y={($(values.north west)-(values.south west)$)}]
      \draw[black,line width=.7pt] (.015,.025) rectangle (.985,.975);
    \end{scope}
  \end{tikzpicture}
  \caption{Actual light-theme Debug observation at \texttt{overflow\_ready}, with decimal register display. The same input pair has overflowed the signed word range but fits in a doubleword. The box is a native LaTeX annotation.}
  \label{fig:revision-fixed-overflow}
\end{figure}


Replace only the probe's \texttt{addu} with \texttt{add} in a disposable file.
Execute it and identify the signed-overflow exception at that instruction.
Explain why it does not contradict the earlier motion range proof: this probe
accepts a different input domain. Restore \texttt{addu} and normal mode before
running the complete program checks.

For multiplication, compute raw $192\cdot320$ and explain its scale before
rescaling. Then use a negative product to distinguish arithmetic right shift
from a logical shift. A shift does not decide the desired rounding rule for you;
state whether you need floor, truncation toward zero, or nearest with a specified
tie convention.

\section{Implement signed interpolation independently}

Open \texttt{fixed-interpolate.asm}. Implement \texttt{interpolate\_raw} with
two sign-extended signed raw words in \texttt{\$a0} and \texttt{\$a1}, and
weight $0\leq w\leq256$ in \texttt{\$a2}. Return the nearest raw word to
$((256-w)a+wb)/256$ in \texttt{\$v0}; exact half ties go toward positive infinity.
Preserve saved registers, the stack pointer, caller-owned stack bytes, inputs
and guards. Temporary registers and private labels are your choice.

Start with the endpoints: weight zero returns $a$ and 256 returns $b$. Predict
the default $(-256,512,64)$ result, raw $-64$. For a rounding boundary, use raw
endpoints $-1$ and zero with weight 128. Their exact midpoint is minus one half
of a raw unit, so the required result is zero. Raw endpoints $-2$ and $-1$ at the
same weight produce $-1$. These are fractions of a raw unit, not half a pixel.

Choose a weighted-sum or difference-based formulation and bound each intermediate.
The wider numerator fits in a signed doubleword, but word products need not fit.
For the difference form, test the most negative and most positive word together:
subtracting them must happen at the wider width too. Delay precision reduction
until the combined numerator is ready, rather than rounding the terms separately.

Use arithmetic that implements floor after the chosen half-unit bias. A signed
division that truncates toward zero gives a different result for some negative
numerators. Verify values immediately around a half tie in both directions.
Then run the public checks:
\begin{verbatim}
java -jar /path/to/cli.jar grade \
  revision/lessons/012-fixed-point/independent/fixed-interpolate.kts \
  /path/to/your/fixed-interpolate.asm
\end{verbatim}
Learn checks the same independent file twice. Its first check has seven
fresh-machine groups containing all 3368 scalar calls from the original public
corpus, followed by an eighth case that runs the caller through its result
store. Each group has at most 512 calls and a budget of 100000 instructions.
The two continuations alternate, global fixture values differ from passed
arguments, and saved registers contain full doubleword patterns. Store traces
reject writes to inputs or the caller's 32 stack bytes even when the bytes are
subsequently restored. Private stack frames below the incoming pointer are
permitted; no instruction sequence is prescribed.

The caller-result case pauses before \texttt{interpolate\_exit}, after the
caller has stored raw $-64$ in \texttt{answer}. It observes the returned
\texttt{\$v0} before the exit selector replaces it. Learn's second check runs
the complete source to silent normal exit with a 1000-instruction budget.
A trap after the result store therefore cannot earn completion. Both checks
must pass for the current source. The command-line script above combines the
scalar corpus with a whole-caller execution check; the JSON profiles packaged
beside it are Learn's data, not executable scripts.

On a mismatch, retain the case name, one-based call number and supplied input
triple from the diagnostic. Derive its exact wide numerator before changing
the rounding or return instruction. An unfinished starter fails explicitly.
A correct default answer does not justify narrower arithmetic for the full
stated domain, and a finite passing corpus does not replace the range proof.

In disposable copies, introduce one word-width product, then one logical shift
of a negative numerator. Find a counterexample for each and explain the first
instruction that loses the needed information. Finally, change the caller's
result store and show why function-only results would miss that error. Restore
the working source and rerun both Learn checks before leaving the exercise.
